\section{Sonuç}
\label{sec:conclusion}

Kullanıcılar hareket ederken bir baz istasyonunun hangi RIS faz optimize edicisini ve ne kadar süreyle koşturması gerektiğini sorduk. Her yöntemin ihtiyaç duyduğu süreyi hesaba katarak şu bulgulara ulaşıyoruz.
\begin{enumerate}
    \item \textbf{Hız, hassasiyeti yener.} Ölçülen adım maliyetlerinde yavaş çözücüler, koordinat taramalarının sayısı ayarlanmış olsa bile, milisaniyeler süren tek kullanıcılı bir sezgisel yönteme yenilir; kanal edinme gecikmesi olmadığında bu, 3 ile 30~km/h arasındaki her hızda geçerlidir (yavaş çözücüler için iyimser değerlerle). Daha uzun pencereler üzerinden ortalandığında koordinat inişi sezgisel yöntemle berabere kalabilir ya da onu geçebilir, ancak ayarlanmış Adam'ın gerisinde kalır.
    \item \textbf{Birkaç adım optimaldir.} Faz birkaç on milisaniyede bir yenilendiğinde Adam'ın dar bant etkin hızı bir avuç adımdan sonra tepe yapar ve ardından düşer (fazı daha uzun tutmak, düz bir optimumu yukarı kaydırır); optimum, 3~km/h hızda 8 adımdan, IMT-2020 yaya sınıfının üst ucu olan 10~km/h hızda 3 adıma kayar.
    \item \textbf{Optimumu tek bir sayı belirler.} En iyi adım sayısını Doppler frekansına göreli adım süresi, $x=\Delta T f_d$, belirler ve bu sayı yaklaşık olarak $x^{-2/(\gamma+2)}$ gibi ölçeklenir (ölçülen üs $-0.675$); yüz kat daha hızlı bir gerçekleştirmeyle yakınsama neredeyse bedava hale gelir.
    \item \textbf{Güncelleme periyodunu pilotlar belirler.} Sürekli işletimde ne sıklıkla güncelleme yapılacağına pilot ek yükü karar verir; kısa güncellemeler (en fazla 8 adım) yeterlidir.
    \item \textbf{Öğrenme saat karşısında kazanmalıdır.} Eşit adımlarda kazanmak yeterli değildir: öğrenilen güncelleme kuralı (statik küme, tek CPU iş parçacığı), her adıma ölçülen süresi maliyet olarak yansıtıldığı anda baskınlanmış olur.
\end{enumerate}
Burada nicelendirilen rejimde gecikme, hareketli RIS için ikincil bir ölçüt değildir: hangi yöntemin en iyi olduğuna o karar verir.

\emph{Pratik kurallar.} Adım sayısını yakınsamaya göre değil; $x$ değerine, kanal edinme gecikmesine ve fazın ne kadar süre tutulacağına göre seçin. Güncelleme periyodunu pilot ek yüküne göre seçin. Yöntemleri, her bütçede ayarlanmış bir klasik taban çizgisine karşı ölçülmüş gecikme üzerinden karşılaştırın ve eşit adımdaki kazanımlar yerine baskınlığı raporlayın.

\emph{Sınırlılıklar.} Ana sonuçlar, BS'nin ön kodlayıcısını güncel kanaldan kusursuz bilgiyle yeniden hesapladığını varsayar. Pilot yayın süresi (\%17--32 daha düşük hızlar), kestirim hatası ve 2--3 bitlik fazlar yavaş çözücüleri ayarlanmış Adam'ın gerisinde bırakır; tek bitte tek taramalı bir arama en iyi yuvarlanmış Adam ile berabere kalır ve 10~dB ve altındaki pilot sinyal-gürültü oranlarında hızlı yöntemler birbirinden ayırt edilemez (Bölüm~\ref{sec:res_robust}). RIS anahtarlama süresi (\cite{Zhao2024OJCOMS} çalışmasında $\ge$0.72~ms) modellenmemiştir ve tek bir geometri vardır. Gecikmeler tek iş parçacıklı Python'dan gelmektedir: yalnızca koordinat inişi derlenirse (62--63 kat daha hızlı) Python Adam'ı geçer, ancak tüm yöntemlerin eşit ölçüde hızlandırılması ayarlanmış Adam'ı önde tutar; normalize adım maliyeti sonuçları başka hızlara ve, yaklaşık olarak, başka taşıyıcılara taşır. Yarı tanımlı gevşetme~\cite{Wu2019TWC}, WMMSE~\cite{Choi2024TCOM} ve manifold yöntemleri zamanlanmamıştır; uzun koşulu Adam ise yalnızca sabit bir bütçede zamanlanmıştır. İlk bir OFDM kontrolünde iç optimum, 20~MHz'de merkez frekans optimizasyonu için geçerliliğini korumuştur; ancak 50--100~MHz'de 48 adım, 16 yapılandırmadan yalnızca 3 tanesinde anlamlı ölçüde daha kötü olmuştur.
